
\section{Evaluation Setup}

\cparagraph{Platforms.}
All experiments are conducted on a GPU workstation equipped with an NVIDIA RTX~3090 GPU (24\,GB memory, 936\,GB/s memory bandwidth) and a dual-socket Intel Xeon Gold 5218R CPU (80 logical cores).
We use PyTorch~2.5.1, Triton~3.1.0, CUDA~12.4, GCC~11.4.0, and Claude Sonnet~4~\cite{anthropic2025claude} as the underlying LLM.

\cparagraph{Benchmarks.}
We evaluate on three benchmark suites:
(1)~\emph{PolyBench/C}~\cite{polybench}: 30 numerical kernels (stencils, linear algebra, data mining) at default sizes ($N$=60--120) and 8$\times$ scale ($N$=480--960);
(2)~\emph{TSVC}~\cite{callahan1988tsvc}: 151 loop kernels covering linear traversals, reductions, indirect addressing, and loop-carried dependences, benchmarked at 32\,MB and 3200\,MB per array;
(3)~8 application kernels from Rodinia~\cite{che2009rodinia} (SRAD, lud, gaussian, lavaMD, hotspot, pathfinder) and ECP proxy applications~\cite{heroux2009mantevo} (miniMD LJ~force, miniFE SpMV).

\cparagraph{Baselines.}
The primary comparison is against \emph{OMP~GPU}---OpenMP GPU offloading compiled by NVIDIA \texttt{nvc} with managed memory, running on the \emph{same} GPU as Triton.
This isolates the quality of LLM-generated Triton code from hardware differences.
We additionally report against \emph{OMP~CPU} (80-thread OpenMP) and \emph{ST} (single-threaded GCC \texttt{-O3}).

\cparagraph{Evaluation methodology.}
All speedups measure kernel-only execution time with data pre-staged on the GPU, averaged over 10--50 iterations after warmup.
For the correctness ablation, we compare \emph{with analysis} (WA) vs.\ \emph{no analysis} (NA, LLM receives only the C source).
For profiling feedback, we report the initial speedup (first correct kernel) and the final speedup after up to 3 NCU-guided optimization iterations.
